\chapter{Wnioski i kierunki dalszych prac}\label{ch:wnioski}

% TODO: podsumowanie — ok. 3--5 stron.

\section{Realizacja celów pracy}\label{sec:cele-zrealizowane}

% TODO: zestawienie celów szczegółowych z~sekcji~\ref{sec:cel-pracy}
% i~ich status (zrealizowane / częściowo / niezrealizowane).

\section{Wkład własny}\label{sec:wklad-wlasny}

% TODO: co jest oryginalnym wkładem autora (vs komponenty open source).

\section{Ograniczenia}\label{sec:ograniczenia}

% TODO: znane ograniczenia: zbiór testowy syntetyczny, brak walidacji na
% prawdziwych danych klinicznych (RODO/IRB), ograniczona liczba specjalizacji,
% zależność od zewnętrznych API LLM w~module chmurowym.

\section{Kierunki dalszego rozwoju}\label{sec:dalsze-prace}

% TODO: walidacja na rzeczywistych danych (po uzyskaniu zgody komisji
% bioetycznej), integracja z~systemami HIS/EHR, rozszerzenie o~kolejne
% specjalizacje, multimodalność (analiza zdjęć, EKG).

\section{Uwagi końcowe}\label{sec:uwagi}

% TODO: refleksja nt. roli AI w~medycynie, odpowiedzialności klinicznej,
% zasady "lekarz ma ostatnie słowo".
